% Conclusion and Future Work

This paper presented AnimatorSBT, a contribution certification framework built on
non-transferable tokens, and used Japanese animation as the setting in which to specify
and cost it. We formalised the contribution invisibility problem as the gap between actual
and credited contributors, derived five requirements from it, and showed why the
institutional route that resolved the same problem in film and television is unavailable
where no body holds craft jurisdiction.

The prototype establishes feasibility and cost structure, not adoption. A soulbound
contract is deployed on a public testnet with an issuance service and a verification
portal, and the studio co-signing layer and a deterministic evidence-hash serialisation
are implemented as tested reference components, 70 tests passing, but are not operated
with independent studio attesters. Two claims therefore remain open, as they have
throughout: verifiability extends to integrity and existence but not to authenticity,
because in a single-operator configuration the issuer also holds the attester role; and
portability and minimal friction rest on wallet abstraction that is designed but not
built.

The cost result generalises. On-chain issuance is not the constraint: a certificate costs
on the order of a tenth of a cent, and that figure is governed less by the contract's logic
than by the length of the metadata pointer it stores. The binding cost is a flat off-chain
pinning subscription, which is also what reintroduces a dependency on a single operator
continuing to pay. Any design that anchors credentials on chain and stores their content
off it inherits this shape, so a per-certificate on-chain cost reported without the storage
arrangement understates the problem.

Three directions follow. The first and most important is empirical: a study with
animators and studios, measuring whether recipients can state what the credential does and
does not prove, and whether studios will co-sign given the transparency it imposes on
them. The gap itself is measurable without any of this machinery, by counting the
contributors a production-management system records against the names its episode's
credits carry---a first estimate of how many people the credits omit, which no published
source provides. The second is to operate the co-signing protocol with independent,
identity-bound attesters, the only thing that converts the authenticity caveat into a
result; the incentive question, why a studio would attest, is the harder half. The third
is to test the framing outside animation, in the other project-based creative industries
where multi-layer subcontracting, capped credits and the disappearance of the producing
organisation recur.

One limit is worth stating without qualification. Kellogg et al.\ survey how
algorithmic instruments introduced into workplaces restructure evaluation and control,
and treat the result as contested terrain in which workers organise against those
instruments rather than being handed power by them~\cite{kellogg2020algorithms}.
Nothing here escapes that, and nothing here substitutes for the organising. A verifiable
record of what an animator did does not change the rate paid for the next cut, the terms on
which it is commissioned, or who decides whether there is more work. It removes one
disadvantage, the inability to show what you have done, and leaves the rest of the
bargaining position where it was. That disadvantage is worth removing on its own terms; we
do not claim removing it is a labour remedy.

The record these certificates would carry already exists. Production generates a
per-cut, per-person trail as a byproduct of getting the work done. What is missing
is not the record but any mechanism that makes it durable, portable and
contestable beyond the organisation that holds it.
